\section{Introduction}

Regular constraints are common in structured prediction: a BIO tag sequence
should be legal, an extracted date should match a format, and a product code may
need to satisfy a finite-state rule.  A constrained decoder can enforce such a
rule at prediction time, but the legality of the decoded output does not answer
how much probability the original model assigned to legal structures.

This paper studies that latter question.  For a CRF posterior
\(\ptheta(y \mid x)\) and a regular language \(\Lang\), we report
\(\ptheta(\Lang \mid x)\), the posterior mass assigned by the original CRF to
the regular-language event.  The denominator remains the original CRF
normalizer over all length-\(T\) label sequences; it is not replaced by a
constrained-support model.

Known product-automaton marginal inference is the computational neighborhood.
Our focus is original-posterior regular-language event mass as an auditable
semantic object, separated from constrained decoding and constrained-support
normalization.  This distinction matters when a legal decoded output coexists
with low posterior mass on the legal event.

The paper makes four conservative claims.  First, posterior regular-language
event mass is a well-defined finite CRF posterior statistic and can be computed
exactly by standard product-state transfer under explicit finite-context local
factorization assumptions.  Second, hard-constrained decoding and posterior
consistency answer different questions.  Third, the event loss
\(-\log \ptheta(\Lang \mid x)\) is a computable training signal, but not an
accuracy theorem.  Fourth, empirical diagnostics show both utility and
boundaries: event risk can rank rule-specific errors, generic uncertainty can
be stronger, event mass can saturate, and misleading rules can harm task
metrics.

\paragraph{Contributions.}
The contribution is not a replacement for WFST, constrained decoding,
regular-constrained CRFs, posterior regularization, or semantic loss.  The
contribution is a posterior-semantics and audit protocol for reporting the
original CRF posterior's mass on regular-language events, with training and
diagnostic experiments that state their limitations.

\subsection{What This Object Answers}

The key distinction is the denominator and the intervention point.  Table
\ref{tab:intro:objects} summarizes the neighboring objects that otherwise look
similar in notation.

\begin{table}[t]
\centering
\small
\caption{Nearby structured-prediction objects answer different questions.}
\label{tab:intro:objects}
\begin{tabular}{llll}
\toprule
Object & Procedure & Space & Question \\
\midrule
CRF posterior &
\(\ptheta(y\mid x)\) &
all \(y\in\Y^T\) &
what distribution did the model assign? \\
Constrained decoding &
\(\arg\max_{y\in\Lang} S_\theta(x,y)\) &
legal outputs &
what is the best legal output? \\
Constrained-support CRF &
\(\ptheta(y\mid x,y\in\Lang)\) &
legal outputs &
what posterior results after support restriction? \\
This audit &
\(\ptheta(\Lang\mid x)=\Z_{\theta,\Lang}(x)/\Z_\theta(x)\) &
all outputs in denominator &
how much original posterior mass is legal? \\
\bottomrule
\end{tabular}
\end{table}

This distinction is useful only if the scalar is reported with its limitations.
If a rule is already saturated, event mass may add little diagnostic
information.  If a rule is weakly related to task correctness, optimizing event
loss can raise event mass while degrading task metrics.  If the goal is only to
return a legal output, a constrained decoder may be the right tool.  The point
of the audit is to make these cases visible rather than to replace those
interventions.

\subsection{Closest-Prior Objection}

A natural objection is that \(\ptheta(\Lang\mid x)\) is just standard marginal
inference on a product automaton with a different name.  We agree about the
computation.  The product-state dynamic program is not the novelty claim.  The
paper's claim is that the resulting original-posterior event probability is a
useful object to report and study: it can be low when constrained Viterbi
returns a legal output, can define a semantic-loss-like training signal, and can
support diagnostics with explicit empirical boundaries.  This is why the
experiments compare posterior event mass, unconstrained decoded legality,
constrained-product decoded legality, task metrics, and uncertainty baselines
as separate quantities.
